% ============================================================
% 分片：s05b-grassmann-intersect  对应原文 PDF 第 58--65 页（印刷页码 58--65）
% 负责章节：5.3 Hypersurfaces in Grassmannians；5.4 Intersections Of Hypersurfaces
% 译者：t01-intro   审校状态：初稿
% 审校：t08-s4 已完成逐式核对（公式修正 0 处；术语 1 处、批注 2 处）
% ============================================================
% 边界：起于第 58 页 "5.3. HYPERSURFACES IN GRASSMANNIANS" 标题，止于第 65 页最后
% 一个完整句（"This model is of interest for two reasons."）。原文第 65 页末句
% "First of all, ... simplest three | generation models ..." 跨至第 66 页顶部，
% 按 Lead 指示由 s05c-topology 收尾，本片不写该句。
% 第 58 页 5.3 标题之前的 "A General Comment" 两段属 §5.2（toric），不在本片范围。
% 记法：原文粗体 CP 按 GLOSSARY/§4.5 写作 $\mathbf{CP}$；WCP 与 G(k,n) 保留原文的
% 粗体写法 $\mathbf{WCP}$、$\mathbf{G}(k,n)$；$W_{eff}$ 按原文作斜体下标。
% 公式以 pages/page-58.png--page-65.png 为唯一依据逐式核对；原文脚注以 \footnote 保留。

% ==== 原文 PDF 第 58 页（印刷页 58）====

\subsection{格拉斯曼流形中的超曲面}

这里给出类似推理的一个简单应用：考虑一个具有非阿贝尔规范群的模型。事实上，我们将找到一个对应：靶空间为格拉斯曼流形（Grassmannian）中的 Calabi--Yau 超曲面的 $\sigma$ 模型，与（这一次规范群为非阿贝尔群的）规范 Landau--Ginzburg 模型相对应。（格拉斯曼流形中的 Calabi--Yau 超曲面在 Hubsch 的书 [41] 中有所讨论。）与先前的经验不同，在这个例子中规范 Landau--Ginzburg 模型的出现将是稳定的：它并不局限于参数的某些特殊取值，%
% ==== 原文 PDF 第 59 页（印刷页 59）====
\label{op:59}
% [跨页接续: 原文第 58 页末 "special values of" 与第 59 页首 "the parameters, and therefore perhaps more interesting." 属同一句，按语义完整保留。]
因而或许更有意思。进一步推广，即采用更复杂的规范群或处于其他表示中的规范群，应当是不言而喻的。

首先，需要给出靶空间为格拉斯曼流形的 $\sigma$ 模型的一种方便实现，例如复 $n$ 维空间中 $k$ 维平面的格拉斯曼流形。为此，从 $kn$ 个手征超场（chiral superfield）$S^{i}{}_{\lambda}$ 出发，$i=1\ldots k$，$\lambda=1\ldots n$。把 $S^{i}{}_{\lambda}$ 视为 $k\times n$ 矩阵 $S$ 的矩阵元；$\bar S$ 表示它的厄米共轭。群 $G=U(k)$ 通过下式作用于诸 $S$：
\begin{equation}\tag{5.20}
S^{i}{}_{\lambda}\to M^{i}{}_{i'}S^{i'}{}_{\lambda},
\end{equation}
其中 $M^{i}{}_{i'}\in U(k)$。为写出规范群为 $U(k)$ 的模型，引入一个 $k\times k$ 的厄米矩阵 $V$，其分量为矢量超场。取拉格朗日量为规范场与物质的动能项，再加上中心因子 $U(1)\subset U(k)$ 的 Fayet--Iliopoulos 项（FI 项）。势能（在 $\sigma=0$ 处）为
\begin{equation}\tag{5.21}
U(s^{i}{}_{\lambda})=\frac{1}{2e^2}\mathrm{Tr}\,D^2
\end{equation}
其中
\begin{equation}\tag{5.22}
-\frac{D}{e^2}=s\bar{s}-r
\end{equation}
（也就是说，$-D/e^2$ 是一个 $k\times k$ 的厄米矩阵，其矩阵元为 $-(D/e^2)^{i}{}_{i'}=\sum_{\lambda}s^{i}{}_{\lambda}\bar s^{\lambda}{}_{i'}-r\delta^{i}{}_{i'}$）。$D$ 的消失可以作如下解释。对固定的 $i$，把诸 $s^{i}{}_{\lambda}$ 视为在 $n$ 维矢量空间 $W\cong\mathbb{C}^n$ 中定义一个矢量 $w^{(i)}$。（如有需要，取 $W$ 的一组固定基 $e^{1}\ldots e^{n}$，并令 $w^{(i)}=\sum_{\lambda}s^{i}{}_{\lambda}e^{\lambda}$。）于是方程 $D=0$ 断言\footnote{对 $r>0$。对 $r<0$，§3.2 中研究的重整化效应对理解模型是重要的；所有真空态都位于大 $\sigma$ 处。}：诸 $w^{(i)}/\sqrt{r}$ 是正交归一的。这些矢量因而张成 $W$ 的一个 $k$ 维子空间 $F$。在除以规范群 $U(k)$ 之后，诸 $s^{i}{}_{\lambda}$ 的期望值%
% ==== 原文 PDF 第 60 页（印刷页 60）====
\label{op:60}
% [跨页接续: 原文第 59 页末 "contained in the expectation" 与第 60 页首 "values of the s^i_λ..." 属同一句，按语义完整保留。]
所含的规范不变信息，无非是 $F$ 的选择。所有 $F$ 构成的空间按定义就是复 $n$ 维空间中 $k$ 维平面的格拉斯曼流形 $\mathbf{G}(k,n)$。因此在低能下，该模型不过是以此为靶空间的 $\sigma$ 模型。用数学语言来说，采用 §4 的说法，上述构造相当于把 $\mathbf{G}(k,n)$ 实现为 $\mathbb{C}^{kn}$ 在 $U(k)$ 作用下的辛商（symplectic quotient）。\footnote{另一种做法是把格拉斯曼流形实现为如下全纯商：允许 $M\in GL(k,\mathbb{C})$，则 (5.20) 定义了 $\mathbb{C}^{kn}$ 上的一个 $GL(k,\mathbb{C})$ 作用；格拉斯曼流形是 $\mathbb{C}^{kn}$ 中一个稠密开集对 $GL(k,\mathbb{C})$ 的商，该开集由包含矩映射零点的好轨道组成。这些恰好就是诸矢量 $w^{(i)}=s^{i}{}_{\lambda}e^{\lambda}$ 线性无关的轨道。}

% [Tutor Note: 推导补足 —— (5.22) 的矩阵形式即 $D=e^2(r\mathbf{1}_k-s\bar s)$，故 $D=0\iff s\bar s=r\mathbf{1}_k$：其分量形式为 $\sum_\lambda s^{i}{}_{\lambda}\bar s^{\lambda}{}_{i'}=r\delta^{i}{}_{i'}$，即诸行矢量 $w^{(i)}$ 两两正交且模方均为 $r$，故 $w^{(i)}/\sqrt{r}$ 构成 $W$ 中的正交归一 $k$ 元组。$r$ 与单位矩阵成正比，只耦合 $D$ 的迹（中心 $U(1)$）部分，$D$ 的无迹部分则给出 $SU(k)$ 的矩映射条件；这正是原脚注“对 $r>0$”的含义——$r<0$ 时该方程无解，须借助 §3.2 的重整化效应理解。把 $D=0$ 的解集模去 $U(k)$ 即得 $\mathbf{G}(k,n)$，这就是下文的辛商。]

% [Tutor Note: 物理图像 —— 上面的辛商与几何不变量理论（GIT）意义下的全纯商是同一件事的两种语言：矩映射零点模去 $U(k)$，与稳定轨道模去 $GL(k,\mathbb{C})$（稳定条件恰为诸 $w^{(i)}$ 线性无关）给出同一个 $\mathbf{G}(k,n)$；这正是 §4 中凯勒商与 GIT 商的对应。在此视角下 FI 参数 $r$ 就是矩映射的水平，$r>0$ 才有零点；而 $[\omega]$ 随 $r$ 线性变化（Duistermaat--Heckman 定理，见 §5.2）。]

为得到该格拉斯曼流形中的一个 Calabi--Yau 超曲面，再引入一个复超场 $P$，它在 (5.20) 下按 $P\to(\det M)^{-n}P$ 变换。然后取超势为 $W=PG(S^{i}{}_{\lambda})$，其中 $G$ 是一个按 $G\to(\det M)^{n}G$ 变换的多项式。$D$ 的公式修正为
\begin{equation}\tag{5.23}
-\frac{D}{e^2}=s\bar{s}-kn|p|^2-r
\end{equation}
标量势中出现的附加项为
\begin{equation}\tag{5.24}
\sum_\alpha|F_\alpha|^2=|G|^2+|p|^2\sum_{i,\lambda}\left|\frac{\partial G}{\partial s^{i}{}_{\lambda}}\right|^2.
\end{equation}

% [Tutor Note: 推导补足与物理图像 —— $P$ 在中心 $U(1)$ 下的荷为 $-kn$（因为 $M=e^{i\alpha}\mathbf{1}_{k}$ 时 $(\det M)^{-n}=e^{-ikn\alpha}$），故 $W=PG$ 的规范不变性要求 $G$ 是 $S$ 的 $kn$ 次齐次多项式；这同时解释了 (5.23) 中 $kn|p|^2$ 的系数：(5.23) 就是在 (5.22) 的 $s\bar s$ 之外再减去 $kn|p|^2$。在 $(2,2)$ 理论中每个手征多重态都含左右行费米子各一，规范反常成对抵消，故 $U(k)$ 与中心 $U(1)$ 的荷分配不受反常约束；这与 §6 的 $(0,2)$ 情形恰成对照。]

对 $r\gg0$，模型可用现已熟悉的方式分析。假设 $G$ 的选取使 $G=0$ 给出的超曲面 $X\subset\mathbf{G}(k,n)$ 光滑，则势的消失意味着 $p=G=0$；低能理论是以 $X$ 为靶空间的 $\sigma$ 模型。对 $r\ll0$，在同样的假设下，模型有唯一的经典真空（模去规范变换），其中 $\langle p\rangle=\sqrt{-r/kn}$，$\langle s^{i}{}_{\lambda}\rangle=0$。$p$ 的期望值把规范群从 $U(k)$ 破缺到 $U(k)$ 中由行列式为一个 $n$ 次单位根的那些矩阵构成的子群 $H$。（$H$ 是 $SU(k)$ 被 $\mathbb{Z}_n$ 扩张得到的群。）低能下的模型是一个规范 Landau--Ginzburg 模型，其无质量手征超场为 $S^{i}{}_{\lambda}$，有效超势为 $W_{eff}=\langle p\rangle G(S^{i}{}_{\lambda})$，规范群为 $H$。

% ==== 原文 PDF 第 61 页（印刷页 61）====
\label{op:61}

\subsubsection*{一个例子}

这里给出一个方便而有趣的例子（见 [41]，第 101 页）。Calabi--Yau 流形 $X$ 的最简单例子之一，是 $\mathbf{CP}^5$ 中由 $G_1=G_2=0$ 给出的完全交（complete intersection），其中 $G_1$、$G_2$ 分别是 $\mathbf{CP}^5$ 的齐次坐标 $T_1,\ldots,T_6$ 的二次与四次齐次多项式。这个模型似乎没有普通的 Landau--Ginzburg 实现。然而它具有一个规范 Landau--Ginzburg 相，因为（对于光滑的二次超曲面 $G_1$）$\mathbf{CP}^5$ 中的轨迹 $G_1=0$ 等价于一份 $\mathbf{G}(2,4)$。因此 $X$ 可以实现为 $\mathbf{G}(2,4)$ 中的一个超曲面，并按上面的方法研究。

这可以如下显式地看出。不要把诸 $T_i$ 想成六分量的“列矢量”，而应把它们排成一个 $4\times4$ 反对称张量的分量：$T_{\lambda\beta}=-T_{\beta\lambda}$，$\lambda,\beta=1\ldots4$。假定 $G_1$ 非退化，则二次多项式 $G_1$ 的选取不含有任何不变信息，因为所有非退化二次型在坐标的线性变换下都等价。$G_1$ 的一个方便取法是
\begin{equation}\tag{5.25}
G_1=\epsilon^{\alpha\beta\gamma\delta}T_{\alpha\beta}T_{\gamma\delta}
\end{equation}
其中 $\epsilon^{\alpha\beta\gamma\delta}$ 是 Levi-Civita 张量。方程 $G_1=0$ 可以用八个复变量 $S^{i}{}_{\lambda}$（$i=1,2$，$\lambda=1\ldots4$）解出：
\begin{equation}\tag{5.26}
T_{\lambda\delta}=\epsilon_{ij}S^{i}{}_{\lambda}S^{j}{}_{\delta}
\end{equation}
其中 $\epsilon_{ij}$ 同样是 Levi-Civita 张量。两个 $S^{i}{}_{\lambda}$ 给出的 $T$ 成比例（因而定义了 $\mathbf{CP}^5$ 中的同一点），当且仅当它们相差 $S^{i}{}_{\lambda}\to M^{i}{}_{i'}S^{i'}{}_{\lambda}$，其中 $M\in GL(2,\mathbb{C})$。由于 $S$ 上的这个等价关系正是导致 $\mathbf{G}(2,4)$ 的那一个，这就完成了 $\mathbf{CP}^5$ 中二次超曲面 $G_1=0$ 与 $\mathbf{G}(2,4)$ 同构的说明。

% [背景补充: (5.26) 中的 $T_{\lambda\delta}$ 正是 $\mathbf{G}(2,4)$ 通过 Plücker 嵌入进入 $\mathbf{CP}(\wedge^{2}\mathbb{C}^{4})=\mathbf{CP}^{5}$ 时的 Plücker 坐标（即 $2\times2$ 子式），而 (5.25) 的二次方程 $G_1=0$ 就是 Plücker 关系（相差归一化），它把 $\mathbf{G}(2,4)$ 截为 $\mathbf{CP}^{5}$ 中的二次超曲面。一般地，$k$ 维平面的格拉斯曼流形通过 Plücker 坐标嵌入 $\mathbf{CP}(\wedge^{k}\mathbb{C}^{n})$。]

我们希望研究由 $\mathbf{CP}^5$ 中方程 $G_1=G_2=0$ 给出的 Calabi--Yau 流形 $X$，其中 $G_1(T_\lambda)$ 如上，$G_2$ 是诸 $T$ 的四次多项式。在代换 (5.26) 之下，$G_2(T)$ 成为 $S$ 的一个八次、$SL(2,\mathbb{C})$ 不变%
% ==== 原文 PDF 第 62 页（印刷页 62）====
\label{op:62}
% [跨页接续: 原文第 61 页末 "SL(2,C) invariant" 与第 62 页首 "polynomial \tilde G(S)." 属同一句，按语义完整保留。]
的多项式 $\tilde G(S)$。$X$ 等价于 $\mathbf{G}(2,4)$ 中的超曲面 $\tilde G(S)=0$。于是，按照我们关于格拉斯曼流形中超曲面的一般讨论，靶空间为 $X$ 的 $\sigma$ 模型可以通过研究具有手征超场 $S^{i}{}_{\lambda}$ 与超势 $W=P\tilde G(S)$ 的 $U(2)$ 规范理论来研究。它的相可以像本节开头那样求出，它们是 Calabi--Yau 相与一个规范 Landau--Ginzburg 相。

\subsection{超曲面的交}

下面简要讨论这些思想对超曲面交的应用。

为避免记号繁杂，只考虑一个最简单的情形：普通射影空间中的超曲面交，例如具有射影坐标 $s_1,\ldots,s_n$ 的 $\mathbf{CP}^{n-1}$。考虑由次数为 $q_a$ 的齐次多项式 $G_a$ 的零点所定义的超曲面 $H_a$，$a=1\ldots k$。诸 $H_a$ 的交 $X$ 是 Calabi--Yau 流形，如果 $\sum_a q_a=n$。$X$ 光滑的条件是诸 $H_a$ 横截（transverse）相交，其含义是：对任何不全为零的复数 $p_a$，方程
\begin{equation}\tag{5.27}
G_a=\sum_a p_a\frac{\partial G_a}{\partial s_i}=0
\end{equation}
只有当 $s_1=\ldots=s_n=0$ 时才有公共解。

% [Tutor Note: 易错预警 —— (5.27) 把两组条件写成了链式等式 $G_a=\sum_a p_a\,\partial G_a/\partial s_i=0$；应读作：对每个 $a$ 有 $G_a=0$，且对每个 $i$ 有 $\sum_a p_a\,\partial G_a/\partial s_i=0$（即诸 $p_a$ 落在 Jacobian 矩阵 $(\partial G_a/\partial s_i)$ 的左零空间中，这正是横截性条件），而不是断言 $G_a$ 等于该线性组合。同样的链式写法也出现在 (5.31)，那里以 $0=G_1=G_2=\cdots$ 开头，读法相同。]

% [Tutor Note: 推导补足 —— 为什么 $\sum_a q_a=n$ 就是 Calabi--Yau 条件：由 $X$ 上的正合列 $0\to T_X\to T_{\mathbf{CP}^{n-1}}|_X\to\bigoplus_a\mathcal{O}(q_a)|_X\to0$ 及 $c_1(T_{\mathbf{CP}^{n-1}})=nH$（$H$ 为超平面类）得 $c_1(T_X)=\big(n-\sum_a q_a\big)H$，故 $c_1(T_X)=0\iff\sum_a q_a=n$。这也解释了 (5.29) 的电荷配置：$s_i$ 的荷取 $1$、$P_a$ 的荷取 $-q_a$，总荷 $n-\sum_a q_a=0$，恰好与第一陈类为零一致。]

为按本文的一般思路描述超曲面 $X$，考虑一个 $U(1)$ 规范理论，其手征超场 $S_i$ 的电荷为 $1$，$P_a$ 的电荷为 $-q_a$。取超势为 $W=\sum_a P_aG_a(S_1,\ldots,S_n)$。普通的势则为
\begin{equation}\tag{5.28}
U=\frac{1}{2e^2}D^2+\sum_\alpha|F_\alpha|^2
\end{equation}
其中
\begin{equation}\tag{5.29}
D=-e^2\left(\sum_i|s_i|^2-\sum_a q_a|p_a|^2-r\right)
\end{equation}
% ==== 原文 PDF 第 63 页（印刷页 63）====
\label{op:63}
以及
\begin{equation}\tag{5.30}
\sum_\alpha|F_\alpha|^2=\sum_a|G_a|^2+\sum_i\left|\sum_a p_a\frac{\partial G_a}{\partial s_i}\right|^2.
\end{equation}

% [Tutor Note: 推导补足 —— (5.30) 由 $W=\sum_a P_aG_a$ 直接求 $F$ 项得到：$\partial W/\partial p_a=G_a$，$\partial W/\partial s_i=\sum_a p_a\,\partial G_a/\partial s_i$，故 $\sum_\alpha|F_\alpha|^2=\sum_a|G_a|^2+\sum_i|\sum_a p_a\,\partial G_a/\partial s_i|^2$。与 (5.24) 对照可见一般模式：每增加一个 $P_a$ 就增加一个 $G_a=0$ 的 $F$ 项条件，这正是完全交在场论中的实现。]

存在两个相。对 $r>0$，利用横截性条件 (5.27) 可以证明（正如我们在此前的类似问题中所做的那样），经典基态的空间是簇 $X$。对 $r<0$，$p$ 有期望值而 $s$ 没有。此时基态的空间是一个加权射影空间（weighted projective space）$W=\mathbf{WCP}^{k-1}_{a_1,\ldots,a_n}$，低能理论是某个以 $W$ 为底空间的矢量丛上的混合 Landau--Ginzburg/$\sigma$ 模型。

% [原文疑误: 此处 $W=\mathbf{WCP}^{k-1}_{a_1,\ldots,a_n}$ 的下标可疑：$p$ 场共有 $k$ 个（$a=1\ldots k$），加权射影空间的权重应由 FI 电荷 $q_a$ 给出，维数 $k-1$ 也对应 $k$ 个权重，故疑应为 $\mathbf{WCP}^{k-1}_{q_1,\ldots,q_k}$；原文写作 $a_1,\ldots,a_n$（$n$ 个下标），与维数和场数不符。此处照抄原样。]

更一般的环面簇（或格拉斯曼流形等）中的超曲面交，可以用这些论述的显然推广来处理。一般而言会得到许多相，涵盖我们见过的所有类型（Calabi--Yau、Landau--Ginzburg、规范 Landau--Ginzburg 以及各种混合型）。人们或许会问：是否存在某个情形，得到的是普通的 Landau--Ginzburg 相，而不是 $\sigma$ 模型/L.G. 混合相？这种情况只在特殊情形下出现，其中一些具有唯象学意义。下面将描述这样一个情形。

\subsubsection*{超曲面交的 Landau--Ginzburg 模型}

在 $\mathbf{CP}^{n-1}\times\mathbf{CP}^{m-1}$ 中（两个因子的齐次坐标分别为 $s_1,\ldots,s_n$ 与 $t_1,\ldots,t_m$），考虑由 $G_1=G_2=0$ 定义的簇 $X$，其中诸 $G$ 是 $s_i$ 与 $t_j$ 的双齐次（bi-homogeneous）多项式。如果 $G_a$（$a=1,2$）关于 $s_i,t_j$ 是次数为 $(q_a,q'_a)$ 的双齐次多项式，则当 $q_1+q_2=n$、$q'_1+q'_2=m$ 时 $X$ 是 Calabi--Yau 流形。$X$ 光滑的条件是：对任何不全为零的复数 $p_a$，方程
\begin{equation}\tag{5.31}
0=G_1=G_2=\sum_a p_a\frac{\partial G_a}{\partial s_i}=\sum_a p_a\frac{\partial G_a}{\partial t_j}
\end{equation}
没有公共解，除非诸 $s_i$ 或诸 $t_j$ 全为零。（使诸 $s_i$ 全为零或诸 $t_j$ 全为零的 $(s_i,t_j)$ 在 $\mathbf{CP}^{n-1}\times\mathbf{CP}^{m-1}$ 中不决定任何点。）然而，这一横截性条件并不足以在我们通常的构造中导致%
% ==== 原文 PDF 第 64 页（印刷页 64）====
\label{op:64}
% [跨页接续: 原文第 63 页末 "not sufficient to lead in our usual construction" 与第 64 页首 "to a Landau-Ginzburg model..." 属同一句，按语义完整保留。]
一个 Landau--Ginzburg 模型；它将导致一个与我们前面见过的一些模型类似的混合 Landau--Ginzburg/$\sigma$ 模型。

为构造一个以 $X$ 的 $\sigma$ 模型为其相之一的模型，考虑 $U(1)\times U(1)$ 规范理论，其手征场为 $S_i$（$i=1\ldots n$，荷 $(1,0)$）、$T_j$（$j=1\ldots m$，荷 $(0,1)$），另有两个场 $P_a$（$a=1,2$，荷 $(-q_a,-q'_a)$）。取超势为 $W=\sum_a P_aG_a(S_i,T_j)$。经典势中相关的部分为
\begin{equation}\tag{5.32}
\begin{aligned}
\sum_a\frac{1}{2e_a^2}D_a^2={}&\frac{e_1^2}{2}\left(\sum_i|s_i|^2-\sum_a q_a|p_a|^2-r_1\right)^2\\
&+\frac{e_2^2}{2}\left(\sum_j|t_j|^2-\sum_a q'_a|p_a|^2-r_2\right)^2
\end{aligned}
\end{equation}
以及
\begin{equation}\tag{5.33}
\begin{aligned}
\sum_\alpha|F_\alpha|^2={}&|G_1|^2+|G_2|^2+\sum_i\left|\sum_a p_a\frac{\partial G_a}{\partial s_i}\right|^2\\
&+\sum_j\left|\sum_a p_a\frac{\partial G_a}{\partial t_j}\right|^2.
\end{aligned}
\end{equation}

% [Tutor Note: 物理图像 —— (5.32) 的 D 项显式展示了电荷矩阵：$S_i$ 的荷为 $(1,0)$、$T_j$ 的荷为 $(0,1)$、$P_a$ 的荷为 $(-q_a,-q'_a)$，故两个矩映射方程分别为 $\sum_i|s_i|^2-\sum_a q_a|p_a|^2=r_1$ 与 $\sum_j|t_j|^2-\sum_a q'_a|p_a|^2=r_2$。$r_1,r_2\gg0$ 时 $P_a$ 不凝聚（$p=0$），$s$ 与 $t$ 分别满足两个射影空间的条件，其交即 $X$；两个 FI 参数可以独立变动，所以相图是二维的。注意 $P_a$（$a=1,2$）的个数等于超曲面方程的个数，而不是规范群因子的个数。]

把这一切设为零，对 $r_1,r_2\gg0$ 可找到一个相，它在低能下描述靶空间为 $X$ 的 $\sigma$ 模型。一般而言存在一个或许颇有意思的多相图，但其中没有一个相描述 Landau--Ginzburg 模型。例如，对 $r_1,r_2\ll0$，令 $D_a=0$ 要求 $p_1,p_2$ 不同时为零。于是诸 $F_\alpha$ 的消失等价于 (5.31)，而横截性条件断言这要求诸 $s_i$ 全为零或诸 $t_j$ 全为零——但不必两者同时为零。若诸 $s_i$ 全为零，则诸 $t_j$ 不被唯一确定，反之亦然。因此人们得到的不是导致 Landau--Ginzburg 模型的孤立经典真空，而是一族真空，并伴有额外的无质量粒子，对应于一个混合模型。这个特定的混合模型实际上具有可约的靶空间（一个分支上 $s_i=0$，另一个分支上 $t_j=0$），并在交 $s_i=t_j=0$ 处有一个奇点（不同于通常的轨形奇点）。

只有当下述条件成立时才会出现 Landau--Ginzburg 模型：$n$、$m$ 以及诸 $q_a,q'_a$ 使得对合适的 $r_1$ 与 $r_2$ 值，经典势的消失迫使诸 $s_i$ 与 $t_j$%
% ==== 原文 PDF 第 65 页（印刷页 65）====
\label{op:65}
% [跨页接续: 原文第 64 页末 "forces the s_i and t_j" 与第 65 页首 "all to vanish." 属同一句，按语义完整保留。]
全部为零。这种情况只在非常特殊的情形下发生。\footnote{看来相关的情形恰好是 $n=m$，$(q_1,q'_1)=(n-1,1)$，$(q_2,q'_2)=(1,n-1)$；或者 $n\ge m$，$(q_1,q'_1)=(n-1,0)$，$(q_2,q'_2)=(1,m)$。}其中一些特殊情形具有唯象学意义。例如取 $n=4$，$m=3$，$(q_1,q'_1)=(3,0)$，$(q_2,q'_2)=(1,3)$。取
\begin{equation}\tag{5.34}
G_1=\sum_{i=1}^{4}S_i^{3},\qquad G_2=\sum_{j=1}^{3}S_iT_i^{3}.
\end{equation}
% [原文疑误: (5.34) 中 $G_2=\sum_{j=1}^{3}S_iT_i^{3}$ 的求和指标 $j$ 与被加项指标 $i$ 不一致；(5.36) 的 $W_{eff}=\sum_aG_a=\sum_{j=1}^{3}\big(S_j^{3}+S_jT_j^{3}\big)+S_4^{3}$ 表明应为 $G_2=\sum_{j=1}^{3}S_jT_j^{3}$。此处照抄原样。]
便有
\begin{equation}\tag{5.35}
\begin{aligned}
\sum_a\frac{1}{2e_a^2}D_a^2={}&\frac{e_1^2}{2}\left(\sum_i|s_i|^2-3|p_1|^2-|p_2|^2-r_1\right)^2\\
&+\frac{e_2^2}{2}\left(\sum_j|t_j|^2-3|p_2|^2-r_2\right)^2
\end{aligned}
\end{equation}
考虑区域 $r_1\ll r_2\ll0$。(5.35) 的消失要求 $p_2\neq0$，并且要么 (a) $p_1\neq0$，要么 (b) $p_1=0$ 且诸 $t_j$ 不全为零。情形 (b) 与 (5.31) 不相容：(5.31) 立刻推出，若 $p_1=0$、$p_2\neq0$，则诸 $t_j$ 全为零。在情形 (a) 中，(5.31) 容易看出蕴含对所有 $i,j$ 有 $s_i=t_j=0$。此时 $|p_1|$ 与 $|p_2|$ 由 (5.35) 唯一确定。因此对这一范围的 $r_1,r_2$，模型有唯一的经典真空（模去规范变换），并在低能下约化为一个 Landau--Ginzburg 轨形（orbifold）。该模型实际上是一个 $\mathbb{Z}_9$ 轨形，因为 $p_1$ 与 $p_2$ 的真空期望值把规范群破缺为 $\mathbb{Z}_9$。无质量超场 $S_i,T_j$ 的有效超势为（通过重新标度变量把 $\langle p_a\rangle$ 取为 1）
\begin{equation}\tag{5.36}
W_{eff}=\sum_a G_a=\sum_{j=1}^{3}\left(S_j^{3}+S_jT_j^{3}\right)+S_4^{3}.
\end{equation}

% [Tutor Note: 推导补足 —— 为什么是 $\mathbb{Z}_9$：$\langle p_1\rangle\neq0$、$\langle p_2\rangle\neq0$ 要求规范变换 $(\alpha,\beta)\in U(1)^2$ 同时保持二者不变，由荷 $(-3,0)$、$(-1,-3)$ 读出 $3\alpha\in2\pi\mathbb{Z}$ 且 $\alpha+3\beta\in2\pi\mathbb{Z}$；$\alpha$ 有 $3$ 个取值，每个 $\alpha$ 下 $\beta$ 有 $3$ 个取值，故剩余规范群有 $9$ 个元素，且由 $(\alpha,\beta)=(2\pi/3,-2\pi/9)$ 生成，确为 $9$ 阶循环群 $\mathbb{Z}_9$。可验证 (5.36) 中每个单项式在此 $\mathbb{Z}_9$ 下都不变，这正是轨形不变超势所要求的。]

这个模型之所以有趣，有两个原因。
% [跨页接续: 本片止于第 65 页末完整句 "This model is of interest for two reasons."；原文下一句 "First of all, ... simplest three | generation models ..." 跨至第 66 页顶部，按 Lead 指示由 s05c-topology 收尾，本片不写该句。]
